\section{The four-summand physical comparison}\label{ext:four-summand-details}
\subsection{A majorant that includes the boundary atom}
\label{ext:closed-half-line}

The Fourier bound below concerns a closed event. Its value at equality
must therefore be one, even when the law has an atom at the threshold.
The following construction proves both this endpoint convention and the
normalization used in the four-summand certificate.

\begin{lemma}[Closed-half-line majorant]\label{ext:prawitz:majorant}
Put
\[
 K(t)=\frac{1-t}{2}+\frac{i}{2}
       \left((1-t)\cot(\pi t)+\frac1\pi\right),\qquad 0<t<1.
\]
For every integrable real random variable $Y$, every real $h$ and every
$A>0$,
\begin{equation}\label{ext:prawitz:expectation}
 \Pp(Y\le h)\le\frac12+
  2\lim_{\epsilon\downarrow0}\Re\int_\epsilon^1
 K(t)e^{-iAth}\E e^{iAtY}\,dt.
\end{equation}
Only the real-part improper limit is intended. The unmodified complex
integral need not converge.
\end{lemma}
\begin{proof}
Define
\[
 W(x)=\left(\frac{\sin\pi x}{\pi x}\right)^2,\quad W(0)=1,
 \qquad
 H_B(x)=\sum_{m\ge1}\{W(x-m)-W(x+m)\}+2xW(x).
\]
The series converges absolutely and locally uniformly, since the shifted
terms are bounded by a constant times $m^{-2}$ on each bounded $x$ interval.
The function $H_B$ is odd, $H_B(0)=0$ and
$H_B(j)=\operatorname{sgn}(j)$ at every nonzero integer.

We first show $|H_B(x)-\operatorname{sgn}(x)|\le W(x)$.
The partial-fraction identity
\[
 \sum_{m\in\mathbb Z}(x-m)^{-2}=\pi^2\csc^2(\pi x)
 \qquad(x\notin\mathbb Z)
\]
follows by integrating $\pi\cot(\pi z)/(z-x)^2$ around the square
with vertices $\pm(N+1/2)\pm i(N+1/2)$. The cotangent is uniformly
bounded on its edges, so the integral is $O(N^{-1})$; the residues are
$(m-x)^{-2}$ at its integer poles and $-\pi^2\csc^2(\pi x)$ at $x$.
For positive noninteger $x$, subtracting this identity from the series
for $H_B$ gives
\[
 H_B(x)-1=\left(\frac{\sin\pi x}{\pi}\right)^2
 \left\{\frac2x-\frac1{x^2}-2\sum_{m\ge1}(x+m)^{-2}\right\}.
\]
Monotonicity of $s\mapsto(x+s)^{-2}$ gives
\[
 \sum_{m\ge1}(x+m)^{-2}\le\frac1x
 \le\frac1{x^2}+\sum_{m\ge1}(x+m)^{-2}.
\]
The braces consequently lie between $-x^{-2}$ and $x^{-2}$.
Integer values, oddness and $W(-x)=W(x)$ prove the claimed bound
on the whole real line. Hence
\[
 M(x):=\frac{1-H_B(x)+W(x)}2\ge\ind_{\{x\le0\}},
 \qquad M(0)=1,
 \qquad 0\le M(x)-\ind_{\{x\le0\}}\le W(x).
\]
In particular $M$ is bounded, and the boundary point is included.

It remains to identify its Fourier kernel. Elementary integration gives
\[
 W(x)=2\int_0^1(1-t)\cos(2\pi tx)\,dt,
 \qquad 2xW(x)=\frac2\pi\int_0^1\sin(2\pi tx)\,dt.
\]
Insert the first identity into the Abel-weighted shifted series with
weight $r^m$, $0<r<1$. Absolute convergence allows termwise integration,
and the difference of the two cosines yields
\[
 \sum_{m\ge1}r^m\{W(x-m)-W(x+m)\}
 =4\int_0^1(1-t)\sin(2\pi tx)
 \frac{r\sin(2\pi t)}{1-2r\cos(2\pi t)+r^2}\,dt.
\]
The fraction tends to $\tfrac12\cot(\pi t)$ and has absolute value
at most $\tfrac12|\cot(\pi t)|$: its denominator equals
$(1-r)^2+4r\sin^2(\pi t)$. The resulting majorant is integrable,
since the sine vanishes at zero and $(1-t)\cot(\pi t)$ is bounded
at one. Dominated convergence therefore gives
\[
 H_B(x)=2\int_0^1
 \left((1-t)\cot(\pi t)+\frac1\pi\right)\sin(2\pi tx)\,dt.
\]
Combining these formulas proves
\[
 M(x)=\frac12+2\lim_{\epsilon\downarrow0}\Re\int_\epsilon^1
                 K(t)e^{2\pi itx}\,dt
     =1+2\Re\int_0^1K(t)(e^{2\pi itx}-1)\,dt.
\]
The second integral is absolutely convergent. For an integrable $U$,
$\E|e^{2\pi itU}-1|\le2\pi t\E|U|$, while $K(t)=O(t^{-1})$
at zero and is bounded near one. Absolute Fubini thus justifies
expectation of this compensated identity. Apply it to
$U=A(Y-h)/(2\pi)$ and use $M(U)\ge\ind_{\{Y\le h\}}$ to obtain
\eqref{ext:prawitz:expectation}.
\end{proof}

For $Z=(X_1+\cdots+X_4)/2$, take $A=6$.
Independence gives $\E e^{6itZ}=\varphi_P(3t)^4$, so
\eqref{ext:prawitz:expectation} has exactly the factor six, raw
frequency $3t$ and positive imaginary kernel used below. Its assumptions
are only integrability and independence; atoms and infinite supports
cause no change to the closed-event bound.
\subsection{Exact elementary enclosures}
\label{ext:elementary-enclosures}

The four-summand certificate uses rational intervals, not floating-point
values as premises. Addition and multiplication take the extrema of
the appropriate endpoint expressions; division first verifies that its
denominator interval excludes zero. At integer scale $S=10^{40}$,
lower and upper endpoints are rounded using floor and ceiling. If
$v=p/q\ge0$, a square-root bracket is checked by the integer inequalities
$qL^2\le pS^2\le qU^2$. Thus the true value stays in the recorded interval
under every permitted arithmetic operation.

The transcendental enclosures have explicit finite remainders. Machin's
identity $\pi=16\arctan(1/5)-4\arctan(1/239)$ and the alternating series
for $\arctan$ give rational endpoints for $\pi$; the first omitted term
bounds each error. The identity follows by the tangent addition formula,
with the angle in $(0,\pi/2)$. Positive square-root and reciprocal
brackets then give $\phi(0)=1/\sqrt{2\pi}$.
For real $z$,
\[
 \left|e^{iz}-\sum_{k=0}^{60}\frac{(iz)^k}{k!}\right|
 \le\frac{|z|^{61}}{61!},
\]
by Taylor's integral remainder and $|e^{is}|=1$.
For $y\ge0$, integration of the Taylor remainder for $e^{-v}$,
$v\ge0$, gives
\[
 \left|\int_0^y e^{-s^2/2}\,ds-
 \sum_{k=0}^{60}\frac{(-1)^k y^{2k+1}}{2^k k!(2k+1)}\right|
 \le\frac{y^{123}}{2^{61}61!\,123}.
\]
Together with $\Phi(-y)=1-\Phi(y)$ this encloses every required normal
CDF value. The proof uses $|d^{61}e^{-v}/dv^{61}|\le1$ on $v\ge0$,
so it does not presume decreasing terms at the start of the series.

Finally, for $0<u<1$ and an integer $m\ge1$,
\[
 \cot(\pi u)=\frac1{\pi u}-\frac{2u}{\pi}
                  \sum_{j\ge1}\frac1{j^2-u^2},\qquad
 \frac1{m+1}\le\sum_{j>m}\frac1{j^2-u^2}
       \le\frac{2m+1}{2m(m+1)}.
\]
The lower tail bound follows by comparison with
$\int_{m+1}^\infty s^{-2}\,ds$; the upper bound follows from the
telescoping sum of $(j^2-1)^{-1}$. The partial-fraction identity is
obtained by integrating the absolutely convergent derivative identity
proved in Lemma~\ref{ext:prawitz:majorant}, with the constant fixed at
$u=1/2$. These formulas enclose every midpoint kernel value.
All strict comparisons used below are comparisons of the resulting
rational endpoints.
% Extracted, normalized physical-class proof from the pinned older manuscript.
% The terminal numerical premise is not discharged by this text alone.
\subsubsection{One defect measure for all frequencies}
For $\eta>0$, let $\nu=W\,dP/\eta$, viewed as a probability on
tagged points $(\sigma,y)$, $\sigma\in\{-1,1\}$, $y>0$.
Write $c_+=b$, $c_-=a$, $p_0=1-w$, and
\[
 k_0(\sigma,y)=\frac1{c_\sigma^2y},\qquad
 k_2(\sigma,y)=\frac1{(y-c_\sigma)^2},\qquad
 k_1(\sigma,y)=\frac1{y-c_\sigma}.
\]
The contacts have zero $\nu$ mass. Singular expressions are integrated
only off the contacts; the following identities prove integrability.

\begin{lemma}\label{ext:n4:lem:common-measure}
With $g(q)=\varphi_Q(q)$ and
$r_q(\sigma,y)=\exp(\mathrm i\sigma q\,\cdot)[0,c_\sigma,c_\sigma,y]$,
\begin{align}
 \varphi_P(q)-g(q)&=\eta\int r_q\,d\nu,\notag\\
 \eta\int k_0\,d\nu&=p_0-P\{0\}\le p_0,\notag\\
 \eta\int_{\sigma}k_2\,d\nu&=u-c_\sigma P\{\sigma c_\sigma\}\le u,
 &\int_{\sigma}k_1\,d\nu&=0.\label{ext:n4:eq:budgets}
\end{align}
The same $\nu$ occurs at every frequency, and $|r_q|\le |q|^3/6$.
For $\eta=0$, $P=Q$ and every characteristic-function difference is
zero, without defining $\nu$.
\end{lemma}
\begin{proof}
Hermite interpolation at $0,c,c$ has remainder
$y(y-c)^2\exp(\mathrm i\sigma q\,\cdot)[0,c,c,y]$.
On each tag, the linear and quadratic terms have the same integrals
against $P$ and $Q$, by their matching first two tagged moments.
Both interpolating polynomials have constant term 1. Their constant
contributions agree only after summing the tags and adding the atom
at zero, since both total masses are 1.
Hermite--Genocchi gives the response bound. Multiplication of the
three costs by $W=y(y-c)^2$ gives
$(y-c)^2/c^2$, $y$, and $y(y-c)$, respectively.
Their integrals give \eqref{ext:n4:eq:budgets}, including the omitted
contact masses. Absolute integrability of $k_1$ follows from Cauchy
and the $k_2$ bound. Lemma~\ref{ext:n4-excess} handles $\eta=0$.
\end{proof}

Fix a threshold $H$ and finite coefficients $\kappa_j(H)$, defined
below. For the bounded class $|X|\le R_*=17/4$, all suprema below
are over $\sigma\in\{-1,1\}$ and $0<y\le R_*$, off the contact.
The linear response is
\[
 A_H(\sigma,y)=\Re\sum_{j=0}^{N-1}
        4\kappa_j(H)g(q_j)^3r_{q_j}(\sigma,y).
\]
For any fixed $\lambda_0,\lambda_+,\lambda_-\ge0$ and
$\theta_+,\theta_-\in\R$, a bound
\begin{equation}\label{ext:n4:eq:priced-supremum}
 S_H\ge\sup_{\sigma,y}
 \{A_H-\lambda_0k_0-\lambda_\sigma k_2-\theta_\sigma k_1\}
\end{equation}
implies
\begin{equation}\label{ext:n4:eq:gamma}
 \eta\int A_H\,d\nu\le
 \Gamma_H(\eta):=p_0\lambda_0+u(\lambda_++\lambda_-)+S_H\eta.
\end{equation}
Upper bounds for $p_0,u$ may replace them. The two equality prices
retain their signs. In particular, $S_H$ and $\Gamma_H$ need not be
nonnegative. The prices are fixed before taking the supremum over
the whole tagged domain. This preserves the cancellation between
frequencies; separate absolute estimates for each $r_q$ do not.

\subsubsection{The fourth-power remainder}
For complex $\psi=g+\delta$,
\begin{equation}\label{ext:n4:eq:quartic}
 \psi^4-g^4=4g^3\delta+
 \delta^2(\psi^2+2\psi g+3g^2).
\end{equation}
Suppose $|g|\le A\le1$, $|\psi|\le1$ and $|\delta|\le E$.
Two complete bounds for the remainder are
\begin{equation}\label{ext:n4:eq:quartic-basic}
 E^2(p^2+2pA+3A^2),\quad p=\min(1,A+E),\qquad
 v(4A^2+2-v),\quad v=\min(E^2,2A^2+1).
\end{equation}
The first follows directly from \eqref{ext:n4:eq:quartic}. For the
second, use
\[
 \psi^4-g^4-4g^3\delta
   =12\delta^2\int_0^1(1-t)(g+t\delta)^2\,dt
\]
and
$|g+t\delta|^2=(1-t)|g|^2+t|\psi|^2-t(1-t)|\delta|^2$.

A sharper two-disk bound is also finite algebra. Rotate $g=a\ge0$
and put $\delta=x+\mathrm i y$, $s=|\delta|^2$. Then
\[
 |\psi^4-g^4-4g^3\delta|^2
 =s^2\{36a^4+4a^2s+s^2+8ax(6a^2+s)+24a^2x^2\}.
\]
For fixed $a,s$, the expression in braces is convex in $x$.
Its maximum on the two disks occurs at a collinear endpoint or on
$|\psi|=1$. Put $r_+=\min(E,1)$, $m=\min(A,1-r_+)$,
$r_-=\min(E,1+A)$. The collinear bounds are
\[
 r_+^2(6m^2+4mr_++r_+^2),\qquad
 r_-^2(6A^2-4Ar_-+r_-^2),\qquad
 \begin{cases}r_-^4,&r_-\le1,\\
 r_-^2(3r_-^2-8r_-+6),&r_-\ge1.
 \end{cases}
\]
On the unit circle it remains to maximize
\begin{equation}\label{ext:n4:eq:lens-polynomial}
 Q_A(s)=s^2\{18A^4+12A^2+6-(12A^2+8)s+3s^2\}
\end{equation}
on $[(1-A)^2,\min(E^2,(1+A)^2)]$. Its endpoints and the sole
interior local maximum
\[
 s_-=(3A^2+2-A\sqrt{4-3A^2})/2
\]
suffice. For $A=0$ the interval is degenerate or empty.
Taking the maximum of the collinear values and the square root of
this last maximum gives the lens bound. To see that no $a<A$ is
lost, note that the unit-circle expression is convex in $a^2$;
its lower feasible endpoint is already collinear. The collinear
expressions increase with their admissible radius after maximization
over $a$. Thus these are whole-disk bounds, not sampled values.
Their minimum with \eqref{ext:n4:eq:quartic-basic} is valid.

On a closed excess interval $[L,U]$, an affine radius
$e_j(\eta)=A_j+B_j\eta$ is used only when it is nonnegative at both
ends and bounds $|\delta_j|$ throughout. Freeze
$p_j=\min(1,a_j+\max\{e_j(L),e_j(U)\})$ in the first bound of
\eqref{ext:n4:eq:quartic-basic}. This gives a quadratic in $\eta$
with nonnegative leading coefficient. Alternatively, freezing $E$
on the whole cell gives a constant disk or lens bound. A choice
between complete bounds is allowed; a coefficientwise minimum is not.

\subsubsection{A physical Prawitz comparison}
We use the closed-CDF Prawitz majorant~\cite{Prawitz1972}, in the
normalization proved in Lemma~\ref{ext:prawitz:majorant}. To distinguish the
factor from the kernel $K$ used there, put
\[
 \widehat K(u)=2K(u)=(1-u)+\mathrm i R(u),\qquad
 R(u)=(1-u)\cot(\pi u)+1/\pi.
\]
For $Z=(X_1+\cdots+X_4)/2$,
\begin{equation}\label{ext:n4:eq:prawitz}
 F_Z(x)\le J_P(x):=\frac12+
 \int_0^1\Re\{\widehat K(u)e^{-6\mathrm i ux}\varphi_P(3u)^4\}\,du.
\end{equation}
The real integrand is regular at zero; no convergence of the complex
integral by itself is asserted. With $D=x-Z$ and $S(u)=uR(u)$ it is
\[
 F_x(u)=(1-u)\E\cos(6Du)+S(u)\E[6D\operatorname{sinc}(6Du)].
\]
The cotangent partial fractions, extended to the endpoints, give
\[
 0\le S\le1/3,\qquad |S'|<2,\qquad |S''|<5.
\]
Here is a direct derivative check. Separate the first pole:
\[
 S=\frac1\pi-\frac2\pi(A+MG),\quad
 A=\frac{u^2}{1+u},\quad M=u^2(1-u),\quad
 G=\sum_{j\ge2}\frac1{j^2-u^2}.
\]
On $[0,1]$, $|A'|\le3/4$, $|A''|\le2$, $|M|\le1/6$,
$|M'|\le1$ and $|M''|\le4$. The derivative series converge
uniformly, and
\[
 G\le3/4,\qquad |G'|\le10/27<1/2,\qquad
 |G''|\le106/135<7/6.
\]
The first sum telescopes. For the other two use
$j^2-u^2\ge3j^2/4$ and
$\sum_{j\ge2}j^{-k}\le2^{-k}+\int_2^\infty t^{-k}\,dt$,
with $k=4,6$. The product rule gives
$|S'|<19/18<2$ and $|S''|<223/54<5$.
For $u>1/2$, put $v=1-u$ and use
$v\cot(\pi v)\le1/\pi$; for $u\le1/2$ the cotangent is
nonnegative. This proves $R\ge0$. The positive partial-fraction
terms then give $0\le S\le1/\pi<1/3$.
Since $\operatorname{sinc}(v)=\int_0^1\cos(tv)\,dt$,
\begin{equation}\label{ext:n4:eq:integrand-curvature}
 |F_x''(u)|\le42\E|D|+108\E D^2+24\E|D|^3.
\end{equation}
These differentiations are justified by the displayed moments.

Now assume $|X|\le R_*=17/4$ and $\rho<46/25$.
The fourth moment follows from these hypotheses; it is not an
additional input. Write $\kappa=\E X^3$, $M_4=\E X^4$.
Lemma~\ref{ext:n4-excess} gives the useful cell bounds
\begin{align}
 |\kappa|&\le\min\{\rho,|b-a|+\eta\},\notag\\
 M_4&\le\min\{R_*\rho, a^2-ab+b^2+
                         (R_*+2\max(a,b))\eta\}.
 \label{ext:n4:eq:local-moments}
\end{align}
For $x=H/2$, choose whole-cell bounds $\kappa_*$ and $M_{4,*}$
from these formulas and put
\[
 D_2=x^2+1,\qquad
 D_4=x^4+6x^2+2|x|\kappa_*+(M_{4,*}+9)/4.
\]
Then $\E|D|\le\sqrt{D_2}$ and
$\E|D|^3\le\sqrt{D_2D_4}$. Uniformly on $|x|\le2$, the simpler
bounds $\E D^2\le5$, $\E D^4<54$, $\E|D|<9/4$,
$\E|D|^3<33/2$ give $|F_x''|<2061/2$.

Take $N=256$, $u_j=(j+1/2)/N$, $q_j=3u_j$ and
\[
 \kappa_j(H)=\widehat K(u_j)e^{-\mathrm i q_jH}/N,\qquad
 C_H(Q)=1+2\Re\sum_j\kappa_j(H)g(q_j)^4-2\Phi(H/2).
\]
Composite midpoint quadrature gives one raw numerator reserve
\begin{equation}\label{ext:n4:eq:midpoint}
 Q_H=\min\left\{\frac{687}{8N^2},
 \frac{42\sqrt{D_2}+108D_2+24\sqrt{D_2D_4}}{12N^2}\right\}.
\end{equation}
It controls the integral of the original law $P$ once. No separate
reference, linear-transfer or remainder quadrature error is added.

The raw thresholds are $H_\ell=-4+\ell/8$, $0\le\ell\le57$.
To fill the intervals between them, bound the characteristic function
on entire frequency cells, not just at the midpoints. If
$a_j\ge|g(q_j)|$ and $e_j$ bounds $|\varphi_P(q_j)-g(q_j)|$ on the
whole excess cell, then
\[
 P_j=\min\{1,a_j+e_j+3/(2N)\}
\]
bounds $|\varphi_P(3u)|$ throughout $j/N\le u\le(j+1)/N$,
since $\E|X|\le1$. Two complete kernel bounds are
\[
 u^2|\widehat K(u)|\le u^2(1-u)+u/3,\qquad
 u^2|\widehat K(u)|\le u^2(1-u)+u(1-u)/2.
\]
The second follows by retaining the first term of the cotangent
partial fractions:
$S(u)\le(1-u)(1+2u)/(\pi(1+u))<(1-u)/2$.
Let $d_j$ be the smaller of the two exact integrals over the $j$th
frequency cell. Their separate totals are $1/4$ and $1/6$.
Thus, for the same whole geometry and excess cell,
\[
 |J_P''(x)|\le B=36\sum_jd_jP_j^4.
\]
Since $\sup|\Phi''|<1/4$, the smooth physical function
$H\mapsto2\{J_P(H/2)-\Phi(H/2)\}$ needs the single interpolation
reserve
\begin{equation}\label{ext:n4:eq:H-reserve}
 E_H=\frac{h^2}{16}(B+1/4),\qquad h=1/8.
\end{equation}
The elementary cube, mass and unit radii
\[
 \eta q_j^3/6,\qquad \max(a,b)^2p_0q_j^2,\qquad 1+a_j
\]
suffice for this curvature calculation. The mass bound follows from
$|r_q|\le q^2/y$ and the first budget in
\eqref{ext:n4:eq:budgets}. To verify the former inequality without
dividing at a contact, put $f(y)=e^{\mathrm i\sigma qy}$. The
divided-difference identity gives
\[
 y f[0,c,c,y]=f[c,c,y]-f[0,c,c].
\]
Each second divided difference has modulus at most $q^2/2$ by
Hermite--Genocchi, including repeated nodes. Hence, for $\eta>0$,
\[
 |\varphi_P(q)-g(q)|
 \le\eta q^2\int\frac1y\,d\nu
 \le q^2\max(a,b)^2p_0.
\]
The zero-excess case has $\varphi_P=g$ and requires no division.

For $H\le-4$, Cantelli gives $F_Z(H/2)\le1/5$.
For $H\ge25/8$, $\Phi(H/2)>4/5$. The latter inequality follows,
for example, by integrating the alternating Gaussian Taylor
polynomial through degree six on $[0,25/16]$ and using
$\phi(0)>19947/50000$. Consequently both threshold tails have
positive normalized score at most $2/5$.

\subsubsection{The low-excess certificate}
On a geometry rectangle and a closed excess cell $[L,U]$, let
$M_H$ bound $C_H(Q)-Tr$ for the entire rectangle. Let
$\Gamma_H$ be obtained from \eqref{ext:n4:eq:gamma}, and let
\[
 R_2(\eta)=r_0+r_1\eta+r_2\eta^2,\qquad r_2\ge0,
\]
bound the sum of the fourth-power remainders weighted by
$|\kappa_j(H)|$. Combining \eqref{ext:n4:eq:quartic} at the same
$H$ gives
\begin{equation}\label{ext:n4:eq:low-polynomial}
 2\{F_{4,P}(H)-\Phi(H/2)\}-T\rho
 \le P_H(\eta):=M_H+Q_H+E_H+2\Gamma_H(\eta)+2R_2(\eta)-T\eta.
\end{equation}
When used at grid points, $E_H$ reserves interpolation to adjacent
thresholds. The leading coefficient is nonnegative. Endpoint
nonpositivity therefore proves the bound on the entire closed excess
cell. Equivalently its degree-two Bernstein coefficients are
\[
 P_H(L),\quad
 \frac{P_H(L)+P_H(U)}2-\frac{p_2(U-L)^2}{2},\quad P_H(U),
\]
where $p_2$ is the quadratic coefficient. For a quantitative signed
score, with $r_L\le r$ and $\rho<46/25$, use
\[
 T+\max_{\eta=L,U}
 \left\{\frac{P_H(\eta)}{46/25},
        \frac{P_H(\eta)}{r_L+\eta}\right\}.
\]
Both denominator faces are needed when the numerator is negative.
The cap branch is convex; for the other branch the derivative
numerator has derivative $P_H''(\eta)(r_L+\eta)\ge0$, so an
interior critical point can only be a minimum.

Here are the finite enclosures that produce these polynomials.
The center is enclosed by signed Taylor expansion in $(w,\beta)$
or in the contacts $(a,b)$ along probability-preserving paths.
For either pair of coordinates $i,j$,
\[
 \partial_i C_H=8\Re\sum_j\kappa_jg^3\partial_i g,\qquad
 |\partial_{ik}C_H|\le8\sum_j|\kappa_j|
       (3D_iD_k+D_{ik}),
\]
where $D_i,D_{ik}$ bound the derivatives of the explicit
three-atom characteristic function. The first-order terms retain
their signs; the second-order error bounds the whole rectangle.
The function $r(w,\beta)$ is jointly convex. Indeed, writing
$r=w^{-1/2}f(\beta)$, $f=(1+\beta^2)/\sqrt{1-\beta^2}$,
the Hessian determinant has the positive factor
\[
 3ff''-(f')^2=
 \frac{9+9\beta^2+15\beta^4-\beta^6}{(1-\beta^2)^3}.
\]
Subtracting its supporting plane therefore encloses $C_H-Tr$
without an absolute Hessian charge for $r$.

For the signed response, put $R_*=17/4$, $y_k=k/64$,
$0\le k\le272$, and use their nearest-node closed cells,
clipped to $[0,R_*]$. Hermite--Genocchi gives
\begin{align}
 r_q(\sigma,y)&=\sum_{m=0}^{60}(\mathrm i\sigma)^{m+3}b_m h_m+E_q,
 &|E_q|&\le\frac{q^3}{6}\frac{(qR_*)^{61}}{61!},\notag\\
 b_0&=q^3/6,&b_m&=b_{m-1}qR_* /(m+3),\notag\\
 h_0&=1,&h_m&=(y/R_*)h_{m-1}+(m+1)(c/R_*)^m.
 \label{ext:n4:eq:response-series}
\end{align}
The bound $|\partial_y r_q|\le q^4/24$ adds the half-cell reserve
$q^4/(48\cdot64)$. The signed frequency coefficients are summed
before bounding their polynomial on a $y$ cell. If $g_0,c_0'$ are
nominal values, $e_q\ge|g-g_0|$ is the geometry error,
$p_q\ge|g|$ and $a_{0,q}\ge|g_0|$, then
\[
 |g^3-g_0^3|\le e_q(p_q^2+p_qa_{0,q}+a_{0,q}^2),
 \qquad |\partial_c r_q|\le q^4/12.
\]
These give a whole-rectangle correction for
$g^3r_q-g_0^3r_{q,0}$, with no second copy of the series or mesh
reserve. Thus the priced supremum bounds actual references throughout
the rectangle, not only its nominal center.

The prices used in the finite certificate are chosen from
\begin{align*}
 \lambda_0&\in\{0,1/1000,1/400,1/200,1/100,1/50,1/25,
                         2/25,4/25,8/25,16/25\},\\
 \lambda_\pm&\in\{0,1/2000,1/1000,1/400,1/200,1/100,1/50\},\\
 \theta_\pm&\in\{0,\pm1/200,\pm1/100,\pm1/50\}.
\end{align*}
For each fixed choice, all 273 cells on each of the two tags enter
the supremum.
When $\lambda_\sigma=0$, only $\theta_\sigma=0$ is used;
otherwise the reciprocal contact rays would give an infinite supremum.
On a cell $[y_L,y_U]$, with $c\in[c_L,c_U]$,
$k_0\ge1/(c_U^2y_U)$. The remaining cost is the concave quadratic
$-\lambda_\sigma z^2-\theta_\sigma z$ in $z=1/(y-c)$.
Its endpoints and admissible vertex suffice; when the cell crosses
the contact both reciprocal rays are included. No division at the
contact or at zero is used for the response series.

The same budgets yield whole-cell radius models by pricing
$|r_q|$ or its signed coordinate projections. They are combined
with the three elementary radii above and the complete two-disk
remainder. An independent absolute branch uses
\[
 |\psi^4-g^4|\le E(p^3+p^2A+pA^2+A^3),\qquad p=\min(1,A+E),
\]
or the complete linear-plus-remainder bound. It is evaluated as a
whole alternative. In every branch, the actual reference modulus,
the actual normalization $\rho=r+\eta$, one midpoint reserve and
one physical threshold reserve are retained. The final maximum over
$H$ is taken only after assembling each complete row.

\subsubsection{Closing the large excess}
Let $p=P(X\ge0)$ and $q=1-p$; zero belongs to the nonnegative class.
Condition on exactly $j$ nonnegative observations among four.
For $Z=(X_1+\cdots+X_4)/2$, the conditional moments are
\[
 \mu_j=\frac{u(j-4p)}{2pq},\qquad
 V_j=\frac14\left[j\left(\frac{v_+}{p}-\frac{u^2}{p^2}\right)
 +(4-j)\left(\frac{v_-}{q}-\frac{u^2}{q^2}\right)\right].
\]
Closed-event Cantelli gives, for every $x$,
\begin{equation}\label{ext:n4:eq:count-cantelli}
 F_Z(x)\le\sum_{j=0}^4\binom4j p^jq^{4-j}A_j(x),\qquad
 A_j(x)=\begin{cases}1,&x\ge\mu_j,\\
 V_j/(V_j+(\mu_j-x)^2),&x<\mu_j.
 \end{cases}
\end{equation}
This includes variance zero and atoms at the threshold. The necessary
sign-mass interval is
\[
 w(1-\beta)/2\le p\le1-w(1+\beta)/2.
\]
It is covered by closed cells. On a threshold cell $[x_L,x_U]$,
bound $A_j(x_U)$ and subtract a lower bound for $\Phi(x_L)$.
The five degree-four Bernstein coefficients bound the binomial
mixture jointly; its masses are not maximized independently.
Each retained sign-mass cell supplies a nonnegative raw numerator
bound $d_i$ and a correlated lower bound $r_i\le r$, so its score
is at most $d_i/(r_i+\eta)$.

If $r_B$ is the whole-box lower bound for $r$, put
\[
 c_C^0=\max\left\{0,\max_i(d_i/T-r_i),\frac{2}{5T}-r_B\right\},
 \qquad
 c_C=10^{-4}\lceil10^4c_C^0\rceil+10^{-4}.
\]
This closes every $\eta\ge c_C$ by a strict bound below $T$,
including the threshold tails. Independently, Shevtsova closes
every $\rho\ge979013/533000$:
\[
 \frac{3031}{10000}\left(1+\frac{323}{500\rho}\right)\le T.
\]
Hence $c_S=\max\{0,979013/533000-r_B\}$ is another closed sufficient
cutoff. Take $c_B=\min(c_C,c_S)$, selecting one complete route.
The low-excess interval $[0,c_B]$ and the large-excess range
$[c_B,\infty)$ include their common endpoint.

\subsubsection{The finite cover and conclusion}
\begin{proposition}[Bounded-law certificate implication]
\label{ext:n4:prop:certificate}
Assume all the terminal inequalities and domain covers in the following
finite scheme have been verified. Then, for every standardized law $P$ with
\[
 |X|\le17/4,\quad \rho(P)<46/25,\quad
 625/2116\le w\le1,\quad |\beta|\le13/20,
\]
and every real $H$,
\[
 \frac{2\{F_{4,P}(H)-\Phi(H/2)\}}{\rho(P)}\le T.
\]
The same assertion holds for the genuine reflected law $-P$.
\end{proposition}
\begin{proof}
Divide $[625/2116,1]$ into 20 equal closed bands and $[0,13/20]$
into 13 closed bands of width $1/20$. There are 260 rectangles.
The minimum of
$r^2=(1+\beta^2)^2/[w(1-\beta^2)]$ is attained at the
largest $w$ and smallest $|\beta|$. A rectangle with
$\min r^2\ge(46/25)^2$ is empty for the strict physical cap.
For every other rectangle construct its own large-excess cutoff.
If $c_B=0$, the large-excess bound covers all $\eta\ge0$ directly.
If $c_B>0$, divide $[0,c_B]$ into 16 equal closed intervals. At each interval
check all 58 thresholds, their 57 joins, the two tails, and both
actual reflected orientations, using the preceding inequalities.
A parent is accepted by its complete bound or by complete bounds
for all four children obtained by bisecting both coordinates.
Each child has its own cutoff; no part of another child's domain
or numerical conclusion is used in its place.

The certificate condition requires a disposition for every one of the
260 closed parent rectangles: a proved necessary-condition contradiction,
a complete direct comparison, or complete comparisons on all four children.
Acceptance requires the complete rational inequalities and their
reconstruction, with no unresolved rectangle.
All elementary operations
are outward rational enclosures; square roots are checked by integer
squaring. Trigonometric, Gaussian and cotangent values are enclosed
by finite series with analytic remainders, using the arithmetic
rules of Section~\ref{ext:elementary-enclosures}. The response uses
\eqref{ext:n4:eq:response-series}, not divided quotients at contacts.
The complete exact replay reconstructs this whole cover and its upper
bounds; its execution and acceptance record are described in
Appendix~\ref{revision:reproducibility}. A sampled value or an optimizer's
status would not supply these inequalities.

For completeness, acceptance has a direct mathematical interpretation.
Every actual canonical point belongs to a closed parent.
In a nonempty parent or child, the large-excess inequalities cover
$\eta\ge c_B$, while the low-excess inequalities cover all of
$[0,c_B]$, including $\eta=0$. The same unknown law supplies the
single $\nu$ in every priced frequency response. Quadrature and
physical interpolation extend the finite rows to every central
threshold; the tails cover the rest. Thus the accepted rational
inequalities prove the proposition for all actual laws, including
bounded infinite supports and zero atoms.
\end{proof}

